\section{Pseudo G-Study Methodology}
\label{app:pseudo-gstudy}

Classical generalizability theory decomposes measurement variation across facets such as persons, items, raters, and occasions. In HAL, the analogous facets are base model, harness, benchmark, and task. A balanced random-effects G-study is not appropriate because the observed response matrix is sparse and highly unbalanced: not every model--harness pair is evaluated on every task.

We initially considered Bernoulli mixed-effects models of the form
\[
\begin{aligned}
y \sim 1
    &+ (1|\text{task}) + (1|\text{model}) + (1|\text{harness})\\
    &+ (1|\text{model}:\text{harness})
     + (1|\text{benchmark}:\text{harness})
     + (1|\text{task}:\text{harness}).
\end{aligned}
\]
This is the closest analogue to a G-study for binary outcomes, but the full Bayesian version was computationally unstable for the current matrix, especially with high-cardinality task--harness effects. We therefore use a staged L2-regularized logistic fixed-effect audit in the main analysis.

For each staged model, we compute Bernoulli log loss,
\[
\ell = -\frac{1}{N}\sum_i \left[y_i\log p_i + (1-y_i)\log(1-p_i)\right],
\]
and deviance,
\[
D = 2N\ell.
\]
We report McFadden-style pseudo-$R^2$,
\[
R^2_{\mathrm{pseudo}} = 1 - \frac{D_{\mathrm{model}}}{D_{\mathrm{null}}},
\]
where the null model is intercept-only. This quantity should be interpreted as relative reduction in Bernoulli deviance, not as ordinary least-squares variance explained. Because all staged models are fit to the same binary rows and likelihood family, the comparisons indicate which facets materially reduce predictive uncertainty.

We do not use a Gaussian linear mixed model as the main audit because the outcome is binary. A linear probability mixed model can produce probabilities outside $[0,1]$, has heteroskedastic residuals by construction, and estimates variance components on a scale that does not match the Bernoulli likelihood used by the IRT model. The staged logistic audit is therefore a more regular and reproducible diagnostic for this response matrix.
